\begin{statementbox}
\textbf{条件：}抛物线 $C:y^2=kx$，$k>0$；圆
$\Gamma:(x-h)^2+y^2=r^2$，圆心 $M(h,0)$、半径 $r>0$。
\textbf{圆心位置必须满足}
\[
\boxed{h=r+\frac{r^2}{k}}.
\]
取抛物线上\textbf{两两不同}的三点 $A_1,A_2,A_3$。这里 $A_iA_j$ 表示两点\textbf{所在的整条直线}。

\begin{center}
% Accurate coordinates: parabola x=y^2, circle center (2,0) radius 1;
% A=(0,0), B=(3,sqrt3), C=(3,-sqrt3).
\begin{tikzpicture}[scale=1.00,>=Stealth,font=\small]
\draw[gray!55,->] (-0.45,0)--(4.45,0) node[right] {$x$};
\draw[gray!55,->] (0,-2.28)--(0,2.28) node[above] {$y$};
\draw[blue!70!black,very thick,domain=-2.05:2.05,samples=120,smooth,variable=\t] plot ({\t*\t},{\t});
\draw[teal!75!black,thick,fill=teal!9] (2,0) circle[radius=1];
\coordinate (A) at (0,0);
\coordinate (B) at (3,{sqrt(3)});
\coordinate (C) at (3,{-sqrt(3)});
\draw[orange!85!black,very thick] (A)--(B)--(C)--cycle;
\coordinate (T1) at (1.5,{sqrt(3)/2});
\coordinate (T2) at (1.5,{-sqrt(3)/2});
\coordinate (T3) at (3,0);
\draw[teal!75!black,densely dashed] (2,0)--(T1) (2,0)--(T2) (2,0)--(T3);
\foreach \p in {A,B,C} \fill[orange!80!black] (\p) circle (1.7pt);
\foreach \p in {T1,T2,T3} \fill[teal!70!black] (\p) circle (1.45pt);
\fill[teal!70!black] (2,0) circle (1.6pt);
\node[below left] at (A) {$A_1=O$};
\node[above right] at (B) {$A_2$};
\node[below right] at (C) {$A_3$};
\node at (2.3,-.4) {$M$};
\node[above right,blue!70!black] at (1,2.2) {$y^2=x$};
\node at (1.35,.35) {\scriptsize $T_1$};
\node at (1.35,-.35) {\scriptsize $T_2$};
\node[right] at (T3) {\scriptsize $T_3$};
\node[right,orange!75!black] at (3.45,-1.05) {\small $A_2A_3:x=3$};
\end{tikzpicture}

\end{center}

\textbf{核心结论：}
\[
\boxed{A_1A_2\text{ 与 }\Gamma\text{ 相切},\quad
A_1A_3\text{ 与 }\Gamma\text{ 相切}
\ \Longrightarrow\ A_2A_3\text{ 与 }\Gamma\text{ 相切}.}
\]

\textbf{速查公式：}设 $P(t)=(t^2/k,t)$，则两不同点的弦、该弦与圆相切的条件分别为
\[
\boxed{P(u)P(v):\ kx-(u+v)y+uv=0},
\]
\[
\boxed{(kh+uv)^2=r^2\bigl[k^2+(u+v)^2\bigr]}.
\]
\textbf{适用边界：}不能省去特殊的圆心条件；竖直弦照样成立；相切是\textbf{直线与圆}的性质，不表示圆必定内切三角形。另一圆心分支 $h=r^2/k-r$ 也有三直线相切的性质，见第 03、05 节。
\end{statementbox}

\clearpage
\begin{statementbox}
\textbf{圆心条件怎样读成位置？}
\begin{center}
% k=r=1: leftmost point L=(r^2/k,0), M=L+(r,0).
\begin{tikzpicture}[x=1.35cm,y=1.35cm,>=Stealth,font=\small]
\draw[gray!55,->] (-.35,0)--(4,0) node[right] {$x$};
\draw[gray!55,->] (0,-1.35)--(0,1.4) node[above] {$y$};
\draw[blue!70!black,thick,domain=-1.15:1.15,samples=80,variable=\t] plot ({\t*\t},\t);
\draw[teal!75!black,thick,fill=teal!9] (2,0) circle(1);
\draw[gray!55,->] (-.35,0)--(4,0);
\foreach \p/\lab in {0/O,1/L,2/M} {\fill (\p,0) circle(1.5pt);\node[above=3pt] at (\p,0) {$\lab$};}
\draw[<->] (0,-1.55)--(1,-1.55) node[midway,below] {$r^2/k$};
\draw[<->] (1,-1.55)--(2,-1.55) node[midway,below] {$r$};
\draw[gray,dotted] (0,-.1)--(0,-1.7) (1,-.1)--(1,-1.7) (2,-.1)--(2,-1.7);
\node[anchor=west,align=left] at (3.35,.9) {$OM=OL+LM$\\[4pt]$h=r^2/k+r$};
\node[align=center] at (2,-2.85) {圆的最左点 $L$ 的横坐标为 $h-r=r^2/k$\\[5pt]图取 $k=r=1$，坐标等比例};
\end{tikzpicture}

\end{center}
先在轴上确定 $L=(r^2/k,0)$，再向右取长度 $r$ 到圆心 $M$，圆的最左点就是 $L$。这给出了主分支 $h=r^2/k+r$ 的位置记法，不能把任意轴上圆都当作符合条件。
\end{statementbox}
